% Figure Legends (250 words per legend)
% - Should always have a short, explanatory title (as well as legend).
% - Legends must explain the figure fully without reference to the text.
% - If applicable, the original source of previously published material must be acknowledged in the Figure legend.
% - Figures must be cited in the main text (box figures must be cited in the box). 

\begin{figure}[p]
  \centering
\includegraphics[width=10cm]{figures/referential-game}
  \caption{\textbf{The referential game of  \citeA{Lazaridou:etal:2017}.} In a referential game, successful communication is the very purpose of the game (as opposed to scenarios in which communication can help players to achieve an independent goal, such as obtaining a valuable object). Referential games have a long history in linguistics, philosophy and game theory \cite{Lewis:1969,Skyrms:2010}. In the game illustrated here, the sender network receives in input two natural images, depicting instances of two distinct categories out of about 500 (here: a dog and a car), with one of the images marked as target (here, the car). The sender processes the images with a convolutional network module and it emits one symbol (sampled from a fixed alphabet), that is given as input to the receiver network, together with the two images (in random order). If the receiver ``points'' to the correct location of the target in the image array (as it does in the figure), both agents are rewarded. The networks are trained by letting them play the game many times, and adjusting their weights based on the reward signal. No supervision is provided about the symbols to be used for communication, so that they are completely free to adapt the emergent protocol to their strategies and biases.}
  \label{fig:referential-game}
\end{figure}

